\documentclass[../main.tex]{subfiles}
\begin{document}

\onlyinsubfile{
  \setcounter{chapter}{1}
  \setcounter{section}{5}
  \setcounter{subsection}{0}
  \addcontentsline{toc}{chapter}{Основы кинематики}
  \addcontentsline{toc}{section}{Задачи преследования}
}

%==============================================================================
\subsection{О чём эта глава}
%==============================================================================

В задачах преследования участвуют два тела: \textit{преследователь}~$A$ и
\textit{цель}~$B$. Цель движется по известному закону, а преследователь выбирает
направление своей скорости по некоторому правилу, связанному с положением цели.
Самое естественное правило~--- в каждый момент бежать прямо на цель; так поступает
собака, догоняющая зайца, или катер, который всё время держит курс на нарушителя.
Такое преследование называют \textit{погоней} (или \textit{чистым преследованием}).
Возможны и другие правила; с двумя из них~--- параллельным сближением и перехватом~---
мы познакомимся в конце главы.

Даже если цель движется совсем просто, например равномерно по прямой, траектория
преследователя оказывается сложной кривой: скорость преследователя всё время
поворачивается вслед за целью. Уравнение такой траектории удаётся найти лишь в
немногих случаях и только с помощью дифференциальных уравнений. Тем не менее
на большинство вопросов о погоне~--- когда преследователь догонит цель, какой путь
он пройдёт, каковы его ускорение и радиус кривизны траектории~--- можно ответить,
не зная траектории. Для этого достаточно двух инструментов, которые у нас уже есть:
\begin{itemize}
  \item переход в систему отсчёта, связанную с одним из тел (раздел об
  относительности движения);
  \item связь нормального ускорения с угловой скоростью вращения вектора скорости,
  $a_n = v\omega_v$, и формула для радиуса кривизны $\rho = \frac{v^2}{a_n}$
  (раздел о криволинейном движении).
\end{itemize}
Главный объект этой главы~--- отрезок, соединяющий преследователя с целью. Мы
научимся описывать, как меняются его длина и направление, и увидим, что почти все
задачи о погоне сводятся к этим двум величинам.

%==============================================================================
\subsection{Линия визирования: скорость сближения и угловая скорость}
%==============================================================================

\begin{definition}
\textit{Линией визирования} называется прямая, соединяющая преследователя~$A$
с целью~$B$. Её положение задают расстоянием $r = |AB|$ и направлением вектора
$\vv{AB}$; единичный вектор этого направления обозначим $\hat{\mbfe}_{AB}$.
\end{definition}

Пусть скорости тел равны $\mbfv_A$ и $\mbfv_B$. Обозначим через $\beta_A$ угол, на
который нужно повернуть вектор $\vv{AB}$ против часовой стрелки, чтобы он стал
сонаправлен со скоростью $\mbfv_A$; угол $\beta_B$ определим так же для скорости
$\mbfv_B$ (рис.~\ref{fig:pur_los}).

\figtikz{images/kin/pur_los.tikz}{Линия визирования $AB$. Углы $\beta_A$ и $\beta_B$ отсчитываются от направления $\vv{AB}$ против часовой стрелки.}{fig:pur_los}

Перейдём в систему отсчёта, которая движется поступательно вместе с
преследователем~$A$. В ней точка~$A$ покоится, а цель движется с относительной
скоростью
\[
\mbfw = \mbfv_B - \mbfv_A.
\]
Разложим $\mbfw$ на две составляющие: \textit{продольную}~$w_\parallel$, направленную
вдоль линии визирования, и \textit{поперечную}~$w_\perp$, перпендикулярную ей.
Проецируя скорости $\mbfv_A$ и $\mbfv_B$ на эти направления, получаем
\[
w_\parallel = v_B\cos\beta_B - v_A\cos\beta_A,
\qquad
w_\perp = v_B\sin\beta_B - v_A\sin\beta_A.
\]
Здесь $w_\perp > 0$, если цель смещается относительно~$A$ против часовой стрелки.

За малое время $\Delta t$ цель смещается относительно~$A$ на $w_\parallel\Delta t$
вдоль линии визирования и на $w_\perp\Delta t$ поперёк неё. По теореме Пифагора новое
расстояние равно
\[
r' = \sqrt{(r + w_\parallel\Delta t)^2 + (w_\perp\Delta t)^2}
\approx r + w_\parallel\Delta t,
\]
поскольку слагаемые порядка $(\Delta t)^2$ пренебрежимо малы по сравнению со
слагаемыми порядка $\Delta t$. Поперечное смещение поворачивает линию визирования на
малый угол
\[
\Delta\theta \approx \frac{w_\perp\Delta t}{r}.
\]
Разделив оба равенства на $\Delta t$, приходим к двум основным формулам главы.

\begin{definition}
\textit{Скорость изменения расстояния} между телами и \textit{угловая скорость
линии визирования} равны
\[
\frac{\Delta r}{\Delta t} = v_B\cos\beta_B - v_A\cos\beta_A,
\qquad
\Omega = \frac{\Delta\theta}{\Delta t} = \frac{v_B\sin\beta_B - v_A\sin\beta_A}{r},
\]
где $\Delta t \to 0$. Величину $v_{\text{сбл}} = -\frac{\Delta r}{\Delta t}$ называют
\textit{скоростью сближения}; при $v_{\text{сбл}} > 0$ тела сближаются. Угловая
скорость $\Omega > 0$, если линия визирования поворачивается против часовой стрелки.
\end{definition}

Обе формулы~--- чистая кинематика относительного движения: они справедливы при
любом движении двух точек, а не только при погоне. По существу это проекции
относительной скорости на линию визирования и на перпендикулярное ей направление.
В компактной векторной записи первая формула выглядит так:
$\frac{\Delta r}{\Delta t} = \hat{\mbfe}_{AB}\cdot(\mbfv_B - \mbfv_A)$.

Полезно запомнить два частных случая. Если цель движется точно вдоль линии
визирования ($\beta_B = 0$ или $\beta_B = \pi$), то вклад цели в поворот линии
визирования равен нулю. Если же цель движется перпендикулярно линии визирования
($\beta_B = \frac{\pi}{2}$), то её вклад в изменение расстояния равен нулю, и вся её
скорость идёт на поворот прямой $AB$.

%==============================================================================
\subsection{Погоня: ускорение и кривизна траектории преследователя}
%==============================================================================

Пусть теперь преследователь~$A$ всё время движется прямо на цель со скоростью,
постоянной по модулю: $v_A = u = \mathrm{const}$, $\beta_A = 0$. Скорость цели
обозначим $v_B = v$, а угол между скоростью цели и линией визирования~--- просто
$\beta = \beta_B$. Основные формулы принимают вид
\[
\frac{\Delta r}{\Delta t} = v\cos\beta - u,
\qquad
\Omega = \frac{v\sin\beta}{r}.
\]

Скорость преследователя всё время направлена вдоль линии визирования, поэтому
вектор $\mbfv_A$ поворачивается вместе с этой прямой: угловая скорость вращения
вектора скорости преследователя равна $\omega_v = \Omega$. Модуль скорости
преследователя не меняется, значит, его тангенциальное ускорение равно нулю, а полное
ускорение совпадает с нормальным:
\[
a_A = a_n = u\,|\Omega| = \frac{uv\,|\sin\beta|}{r}.
\]
Радиус кривизны траектории преследователя находим по формуле
$\rho = \frac{u^2}{a_n}$:
\[
\rho_A = \frac{u}{|\Omega|} = \frac{ur}{v\,|\sin\beta|}.
\]
Обратите внимание: при сближении ($r \to 0$) ускорение преследователя растёт, а
радиус кривизны уменьшается, если только цель не оказывается строго впереди или
позади него ($\sin\beta \to 0$). Поэтому в последние мгновения погони траектория
преследователя, как правило, закручивается всё круче.

\begin{example}
Лиса бежит равномерно и прямолинейно со скоростью $v$. Собака преследует её со
скоростью $u$, постоянной по модулю и всё время направленной на лису. В некоторый
момент скорости собаки и лисы перпендикулярны, а расстояние между ними равно $L$.
Найдите ускорение собаки и радиус кривизны её траектории в этот момент.

Скорость собаки направлена вдоль линии визирования, а скорость лисы ей
перпендикулярна, поэтому $\beta = \frac{\pi}{2}$. Лиса целиком <<поперечна>>: её
скорость поворачивает линию визирования с угловой скоростью
\[
\Omega = \frac{v}{L}.
\]
С той же угловой скоростью поворачивается вектор скорости собаки, поэтому
\[
a = u\Omega = \frac{uv}{L},
\qquad
\rho = \frac{u^2}{a} = \frac{uL}{v}.
\]
Отметим, что в этот момент расстояние между телами уменьшается со скоростью
$v_{\text{сбл}} = u - v\cos\frac{\pi}{2} = u$, как если бы лиса стояла на месте.
\end{example}

\begin{example}
Преследователь движется со скоростью $u$ прямо на цель, которая движется со
скоростью $v$ по окружности радиуса $R$ вокруг преследователя. Может ли
преследователь приблизиться к цели?

Цель всё время движется перпендикулярно линии визирования: $\beta = \frac{\pi}{2}$.
В начальный момент $\frac{\Delta r}{\Delta t} = v\cos\frac{\pi}{2} - u = -u$, т.\,е.
расстояние сокращается так же, как при неподвижной цели. Движение цели
перпендикулярно линии визирования не мешает сближению~--- оно только поворачивает
эту линию, заставляя преследователя бежать по кривой. Разумеется, как только
преследователь сместится из центра окружности, угол $\beta$ перестанет быть прямым,
и задача станет сложнее; её мы разберём в конце главы.
\end{example}

%==============================================================================
\subsection{Симметричная погоня: черепахи в вершинах многоугольника}
%==============================================================================

Самые красивые задачи о погоне~--- симметричные. Классическая задача: в вершинах
правильного $n$-угольника со стороной $a$ сидят $n$ черепах. Одновременно все они
начинают ползти с одинаковыми по модулю скоростями $v$, каждая~--- всё время прямо
на соседку, следующую за ней по часовой стрелке. Через какое время черепахи
встретятся и какой путь проползёт каждая?

\noindent\textit{Симметрия.} Повернём мысленно всю картину на угол $\frac{2\pi}{n}$
вокруг центра $O$ многоугольника. Каждая черепаха займёт место соседки, и новое
расположение ничем не будет отличаться от исходного: условия движения у всех черепах
одинаковы. Поэтому в любой момент времени черепахи остаются в вершинах правильного
$n$-угольника с тем же центром~$O$; многоугольник лишь поворачивается и сжимается.
Встретятся черепахи в точке~$O$.

\noindent\textit{Время встречи.} Пусть $A$ и $B$~--- соседние черепахи, $r = |AB|$~---
сторона текущего многоугольника. Черепаха~$A$ ползёт прямо на~$B$, т.\,е. вдоль
стороны $AB$. Черепаха~$B$ ползёт вдоль следующей стороны $BC$, которая составляет с
продолжением стороны $AB$ внешний угол правильного $n$-угольника, $\beta = \frac{2\pi}{n}$.
Следовательно,
\[
v_{\text{сбл}} = u - v\cos\beta = v\left(1 - \cos\frac{2\pi}{n}\right) = 2v\sin^2\frac{\pi}{n}.
\]
Скорость сближения постоянна, поэтому сторона многоугольника убывает равномерно,
$r(t) = a - v_{\text{сбл}}t$, и черепахи встретятся через время
\[
T = \frac{a}{v_{\text{сбл}}} = \frac{a}{2v\sin^2\frac{\pi}{n}}.
\]
Модуль скорости каждой черепахи постоянен, поэтому её путь до встречи
\[
s = vT = \frac{a}{2\sin^2\frac{\pi}{n}}.
\]

\begin{center}
\begin{tabular}{c|c|c|c}
$n$ & $\beta$ & $T$ & $s$ \\ \hline
$3$ & $120^\circ$ & $\frac{2a}{3v}$ & $\frac{2a}{3}$ \\[4pt]
$4$ & $90^\circ$  & $\frac{a}{v}$    & $a$ \\[4pt]
$6$ & $60^\circ$  & $\frac{2a}{v}$   & $2a$
\end{tabular}
\end{center}

Тот же результат получается, если следить не за стороной, а за расстоянием $R$ от
черепахи до центра. Радиус описанной окружности $R_0 = \frac{a}{2\sin\frac{\pi}{n}}$,
а угол между стороной многоугольника и радиусом, проведённым в вершину, равен
$\frac{\pi}{2} - \frac{\pi}{n}$. Поэтому скорость черепахи раскладывается на
составляющую, направленную к центру, $v\sin\frac{\pi}{n}$, и составляющую,
перпендикулярную радиусу, $v\cos\frac{\pi}{n}$. Первая постоянна, и
\[
T = \frac{R_0}{v\sin\frac{\pi}{n}} = \frac{a}{2v\sin^2\frac{\pi}{n}}.
\]

\noindent\textit{Ускорение и радиус кривизны.} Линия визирования $AB$ поворачивается
с угловой скоростью
\[
\Omega = \frac{v\sin\beta}{r} = \frac{v\sin\frac{2\pi}{n}}{r},
\]
поэтому ускорение черепахи и радиус кривизны её траектории равны
\[
a_n = v\Omega = \frac{v^2\sin\frac{2\pi}{n}}{r},
\qquad
\rho = \frac{r}{\sin\frac{2\pi}{n}}.
\]
Для квадрата ($n = 4$) отсюда $\rho = r$: в начальный момент $\rho_0 = a$, а через
время $\frac{T}{2}$, когда $r = \frac{a}{2}$, радиус кривизны уменьшается вдвое, а
ускорение вдвое растёт.

\figtikz{images/kin/pur_turtles4.tikz}{Траектории четырёх черепах~--- логарифмические спирали, закручивающиеся к центру квадрата.}{fig:pur_turtles4}

\noindent\textit{Форма траектории.} Скорость черепахи составляет с направлением на
центр постоянный угол $\psi = \frac{\pi}{2} - \frac{\pi}{n}$. Кривая, касательная к
которой образует постоянный угол с радиус-вектором, проведённым из фиксированной
точки, называется \textit{логарифмической спиралью} (рис.~\ref{fig:pur_turtles4}).
За малое время $\Delta t$ черепаха приближается к центру на
$\Delta R = -v\sin\frac{\pi}{n}\,\Delta t$ и поворачивается вокруг него на угол
$\Delta\varphi = \frac{v\cos\frac{\pi}{n}}{R}\,\Delta t$, так что
\[
\Delta\varphi = -\ctg\frac{\pi}{n}\cdot\frac{\Delta R}{R}.
\]
Угол поворота определяется \emph{относительным} изменением расстояния до центра.
Каждый раз, когда расстояние до центра уменьшается вдвое, черепаха поворачивается на
один и тот же угол. Уменьшать расстояние вдвое можно бесконечно много раз, поэтому
до встречи каждая черепаха совершит бесконечно много оборотов вокруг центра.
При этом путь её конечен и равен $s = \frac{a}{2\sin^2\frac{\pi}{n}}$! Для знакомых
с логарифмом: суммирование равенства даёт
$\varphi = \ctg\frac{\pi}{n}\,\ln\frac{R_0}{R}$, или $R = R_0e^{-\varphi\tg(\pi/n)}$.

Радиус кривизны логарифмической спирали равен $\rho = \frac{R}{\sin\psi}$, где $R$~---
расстояние до центра. Проверим согласие с полученной выше формулой: $\sin\psi =
\cos\frac{\pi}{n}$, а сторона многоугольника $r = 2R\sin\frac{\pi}{n}$, поэтому
\[
\frac{r}{\sin\frac{2\pi}{n}} = \frac{2R\sin\frac{\pi}{n}}{2\sin\frac{\pi}{n}\cos\frac{\pi}{n}}
= \frac{R}{\cos\frac{\pi}{n}}.
\]
Две совершенно разные цепочки рассуждений привели к одному и тому же результату.

%==============================================================================
\subsection{Цель движется по прямой: инвариант погони}
%==============================================================================

Рассмотрим основную задачу теории погони. Цель~$B$ движется равномерно и
прямолинейно со скоростью $v$. Преследователь~$A$ движется со скоростью $u$,
постоянной по модулю и всё время направленной на цель. Обозначим через $\varphi$
угол между линией визирования $\vv{AB}$ и скоростью цели, а через
\[
x = r\cos\varphi
\]
проекцию отрезка $AB$ на направление движения цели~--- на сколько цель <<опережает>>
преследователя вдоль своей прямой.

За малое время $\Delta t$ расстояние между телами меняется на
\[
\Delta r = (v\cos\varphi - u)\Delta t.
\]
Проекция $x$ меняется за счёт движения цели на $v\Delta t$ и за счёт движения
преследователя на $-u\cos\varphi\,\Delta t$ (скорость преследователя направлена
вдоль $AB$, и её проекция на направление движения цели равна $u\cos\varphi$):
\[
\Delta x = (v - u\cos\varphi)\Delta t.
\]
Умножим первое равенство на $u$, второе~--- на $v$ и сложим. Слагаемые, содержащие
неизвестный угол $\varphi$, сокращаются:
\[
u\,\Delta r + v\,\Delta x = (v^2 - u^2)\Delta t.
\]
Величины $u$ и $v$ постоянны, поэтому слева стоит изменение комбинации $ur + vx$.
Она меняется равномерно, с постоянной скоростью $v^2 - u^2$.

\begin{theorem}
\textbf{Инвариант погони.} Если цель движется равномерно и прямолинейно со
скоростью $v$, а преследователь бежит прямо на неё с постоянной по модулю скоростью
$u$, то
\[
ur + vx = ur_0 + vx_0 - (u^2 - v^2)t,
\]
где $r$~--- расстояние между телами, $x$~--- проекция отрезка $AB$ на направление
движения цели, индекс $0$ отмечает начальные значения.
\end{theorem}

Эта теорема позволяет ответить на основные вопросы о погоне, не зная траектории
преследователя. Разберём три случая.

\noindent\textit{1. Преследователь быстрее цели ($u > v$).} Комбинация $ur + vx$
убывает и в момент поимки обращается в нуль, так как $r = x = 0$. Время погони
\[
T = \frac{ur_0 + vx_0}{u^2 - v^2}.
\]
Если в начальный момент скорость цели перпендикулярна линии визирования
($x_0 = 0$, $r_0 = L$), то
\[
T = \frac{uL}{u^2 - v^2},
\qquad
s_A = uT = \frac{u^2L}{u^2 - v^2},
\qquad
s_B = vT = \frac{uvL}{u^2 - v^2},
\]
где $s_A$ и $s_B$~--- пути преследователя и цели до встречи.

\noindent\textit{2. Скорости равны ($u = v$).} Тогда $r + x = r_0 + x_0 = \mathrm{const}$.
Поскольку $x = r\cos\varphi$, имеем $r(1 + \cos\varphi) = r_0 + x_0$. Расстояние
всё время уменьшается, так как $\frac{\Delta r}{\Delta t} = -v(1 - \cos\varphi) < 0$
при $\varphi \neq 0$, но не может стать меньше $\frac{r_0 + x_0}{2}$, потому что
$1 + \cos\varphi \leq 2$. Если бы угол $\varphi$ всё время оставался заметно
отличным от нуля, расстояние убывало бы с конечной скоростью и рано или поздно
нарушило бы это ограничение. Следовательно, $\varphi \to 0$: собака постепенно
выходит точно в хвост лисе, а расстояние между ними стремится к пределу
\[
r_\infty = \frac{r_0 + x_0}{2}.
\]
При перпендикулярном начальном положении ($x_0 = 0$) собака так и не догонит лису,
а расстояние между ними стремится к $\frac{L}{2}$.

\noindent\textit{3. Преследователь медленнее цели ($u < v$).} Комбинация $ur + vx$
растёт неограниченно, и догнать цель невозможно.

\begin{example}
На берегу находятся пункты $A$ и $B$, удалённые на расстояние $L$ друг от друга.
Контрабандисты отправляются из пункта $A$ в открытое море перпендикулярно берегу с
постоянной скоростью $v$. Офицер береговой охраны, обнаружив нарушителей в момент
отчаливания их судна, тотчас пускается в погоню на катере из точки $B$. Катер всё
время держит курс на нарушителей и, двигаясь с постоянной скоростью $u$, настигает
судно на расстоянии $L$ от берега. Во сколько раз $k = \frac{u}{v}$ скорость катера
превосходит скорость судна?

В начальный момент линия визирования (от катера к судну) идёт вдоль берега, а
судно движется перпендикулярно берегу, т.\,е. перпендикулярно линии визирования. По инварианту
погони время погони $T = \frac{uL}{u^2 - v^2}$. С другой стороны, судно удалилось
от берега на $L$, поэтому $T = \frac{L}{v}$. Приравнивая, получаем
\[
uv = u^2 - v^2
\quad\Rightarrow\quad
k^2 - k - 1 = 0
\quad\Rightarrow\quad
k = \frac{1 + \sqrt{5}}{2} \approx 1{,}618.
\]
Отношение скоростей равно золотому сечению.
\end{example}

\begin{example}
Собака и лиса бегут с одинаковыми скоростями $v$; лиса~--- по прямой, собака~--- всё
время прямо на лису. В начальный момент расстояние между ними равно $L$, а скорость
лисы составляет с линией визирования угол $60^\circ$ (лиса убегает). Какое расстояние
установится между ними?

В начальный момент $x_0 = L\cos 60^\circ = \frac{L}{2}$, поэтому
\[
r_\infty = \frac{r_0 + x_0}{2} = \frac{L + \frac{L}{2}}{2} = \frac{3L}{4}.
\]
Если бы лиса в начальный момент бежала навстречу собаке под тем же углом
($x_0 = -\frac{L}{2}$), установившееся расстояние было бы равно $\frac{L}{4}$.
\end{example}

%==============================================================================
\subsection{Уравнение кривой погони}
%==============================================================================

\textit{Этот подраздел рассчитан на читателей, знакомых с производной и
интегралом; его можно пропустить без ущерба для понимания остального.}

Найдём траекторию преследователя в задаче о цели, движущейся по прямой (эта задача
была решена П.~Бугером в 1732~г.). Пусть цель стартует из начала координат и
движется вдоль оси $Oy$ со скоростью $v$, а преследователь стартует из точки $(L, 0)$
и движется со скоростью $u$. Обозначим $\lambda = \frac{v}{u}$. Траекторию
преследователя будем искать в виде $y = y(x)$, где $x$ убывает от $L$ до $0$.

В момент $t$ цель находится в точке $(0,\, vt)$. Скорость преследователя направлена
на цель, т.\,е. касательная к траектории в точке $(x, y)$ проходит через точку
$(0,\, vt)$. Уравнение касательной $Y = y + y'(X - x)$ при $X = 0$ даёт
\[
y - xy' = vt = \lambda s,
\]
где $s = ut$~--- путь, пройденный преследователем. Продифференцируем это равенство
по $x$. Слева получаем $y' - y' - xy'' = -xy''$. Путь растёт при убывании $x$,
поэтому $\frac{\dd{s}}{\dd{x}} = -\sqrt{1 + y'^2}$. Следовательно,
\[
xy'' = \lambda\sqrt{1 + y'^2}.
\]
Это дифференциальное уравнение первого порядка относительно $p = y'$. Разделяя
переменные, получаем
\[
\frac{\dd{p}}{\sqrt{1 + p^2}} = \lambda\frac{\dd{x}}{x}.
\]
В начальный момент преследователь бежит к началу координат вдоль оси $Ox$, поэтому
$p(L) = 0$. Интегрируя от $L$ до $x$, находим
\[
\ln\left(p + \sqrt{1 + p^2}\right) = \lambda\ln\frac{x}{L},
\qquad
p + \sqrt{1 + p^2} = \left(\frac{x}{L}\right)^{\lambda}.
\]
Отсюда (выразив $\sqrt{1 + p^2}$ и возведя в квадрат)
\[
y' = p = \frac{1}{2}\left[\left(\frac{x}{L}\right)^{\lambda} - \left(\frac{x}{L}\right)^{-\lambda}\right].
\]
Ещё одно интегрирование с условием $y(L) = 0$ при $\lambda \neq 1$ даёт уравнение
кривой погони:
\[
y = \frac{L}{2}\left[\frac{1}{1 + \lambda}\left(\frac{x}{L}\right)^{1 + \lambda}
- \frac{1}{1 - \lambda}\left(\frac{x}{L}\right)^{1 - \lambda}\right]
+ \frac{\lambda L}{1 - \lambda^2}.
\]

\figtikz{images/kin/pur_curve.tikz}{Кривые погони при $\lambda = \frac{1}{2}$ (сплошная линия; преследователь догоняет цель в точке $C$) и при $\lambda = 1$ (штриховая линия; преследователь не догоняет цель). Касательная к траектории в точке $A$ проходит через текущее положение цели $B$.}{fig:pur_curve}

При $\lambda < 1$ и $x \to 0$ первые два слагаемых стремятся к нулю, и
преследователь догоняет цель в точке
\[
y_C = \frac{\lambda L}{1 - \lambda^2} = \frac{uvL}{u^2 - v^2},
\]
что в точности совпадает с путём цели $s_B$, найденным из инварианта погони.

При $\lambda = 1$ интегрирование даёт логарифм:
\[
y = \frac{x^2}{4L} - \frac{L}{2}\ln\frac{x}{L} - \frac{L}{4}.
\]
При $x \to 0$ ордината неограниченно растёт: преследователь не догоняет цель.
Расстояние до цели найдём из того, что цель лежит на касательной:
$r = \sqrt{x^2 + (xy')^2} = x\sqrt{1 + p^2}$. При $\lambda = 1$ имеем
$\sqrt{1 + p^2} = \frac{1}{2}\left(\frac{x}{L} + \frac{L}{x}\right)$, так что
\[
r = \frac{x^2 + L^2}{2L} \to \frac{L}{2} \quad\text{при } x \to 0.
\]
Получили тот же предел $\frac{L}{2}$, что и из инварианта погони.

Наконец, проверим формулу для радиуса кривизны. В начальной точке $x = L$ имеем
$p = 0$ и $y'' = \frac{\lambda}{L}$, поэтому
\[
\rho = \frac{(1 + y'^2)^{3/2}}{|y''|} = \frac{L}{\lambda} = \frac{uL}{v},
\]
как и в примере с собакой и лисой: в начальный момент скорость цели
перпендикулярна линии визирования.

%==============================================================================
\subsection{Другие стратегии: параллельное сближение и перехват}
%==============================================================================

Погоня~--- не самый выгодный способ догнать цель: преследователь всё время бежит
туда, где цель \emph{находится}, а не туда, где она \emph{будет}. Рассмотрим другое
правило.

\begin{definition}
\textit{Параллельным сближением} называется такое движение преследователя, при
котором линия визирования не поворачивается: $\Omega = 0$, т.\,е.
\[
u\sin\beta_A = v\sin\beta_B,
\]
и при этом расстояние между телами уменьшается:
$v_{\text{сбл}} = u\cos\beta_A - v\cos\beta_B > 0$.
\end{definition}

В системе отсчёта цели преследователь при параллельном сближении движется точно
по линии визирования, т.\,е. по прямой. Если цель движется равномерно и
прямолинейно, а модуль скорости преследователя постоянен, то постоянен и угол
$\beta_A$: скорость преследователя не меняется, и он движется по прямой прямо в
точку будущей встречи (рис.~\ref{fig:pur_intercept}). Такое движение называют
\textit{перехватом}. Линия визирования при этом переносится параллельно самой себе,
отсюда и название метода. Моряки знают это правило давно: если пеленг на
встречное судно не меняется, а расстояние сокращается, суда столкнутся.

\figtikz{images/kin/pur_intercept.tikz}{Перехват. Преследователь движется по прямой $AM$, линия визирования ($AB$, затем $A'B'$) переносится параллельно самой себе.}{fig:pur_intercept}

Из условия $u\sin\beta_A = v\sin\beta_B$ видно, что перехват возможен не всегда.
Если $u > v$, подходящий угол $\beta_A$ найдётся при любом $\beta_B$, и сближение
обеспечено. Если же $u \leq v$, необходимо, чтобы $u \geq v|\sin\beta_B|$ и чтобы
цель двигалась навстречу, т.\,е. $\cos\beta_B < 0$.

\begin{example}
Судно идёт прямым курсом со скоростью $v$. В начальный момент катер находится на
расстоянии $L$ от судна, причём отрезок, проведённый от судна к катеру, составляет с
курсом судна острый угол $\alpha$. С какой наименьшей скоростью должен двигаться
катер, чтобы встретиться с судном? Через какое время произойдёт встреча?

Перейдём в систему отсчёта судна. В ней катер должен двигаться по прямой прямо на
судно, т.\,е. его скорость относительно судна $\mbfu - \mbfv$ должна быть направлена
вдоль линии визирования от катера к судну. Значит, конец вектора $\mbfu$ лежит на
прямой, проведённой через конец вектора $\mbfv$ параллельно линии визирования.
Наименьший вектор $\mbfu$~--- перпендикуляр, опущенный на эту прямую; его длина равна
составляющей скорости судна, перпендикулярной линии визирования:
\[
u_{\min} = v\sin\alpha.
\]
Катер при этом движется перпендикулярно линии визирования, и скорость сближения
равна составляющей скорости судна вдоль линии визирования, $v\cos\alpha$. Время до
встречи
\[
T = \frac{L}{v\cos\alpha}.
\]
\end{example}

Сравним погоню и перехват в простейшей ситуации: цель движется со скоростью $v$
перпендикулярно начальной линии визирования, преследователь (скорость $u > v$)
находится на расстоянии $L$. При перехвате преследователь бежит по гипотенузе
прямоугольного треугольника с катетами $L$ и $vT$, так что $(uT)^2 = L^2 + (vT)^2$ и
\[
T_{\text{пер}} = \frac{L}{\sqrt{u^2 - v^2}}.
\]
При погоне $T_{\text{пог}} = \frac{uL}{u^2 - v^2}$, так что
\[
\frac{T_{\text{пог}}}{T_{\text{пер}}} = \frac{u}{\sqrt{u^2 - v^2}} > 1.
\]
Например, при $u = 2v$ погоня длится $\frac{2L}{3v}$, а перехват~--- всего
$\frac{L}{\sqrt{3}\,v} \approx 0{,}58\,\frac{L}{v}$. Чем ближе скорость
преследователя к скорости цели, тем сильнее проигрыш погони.

Полезно знать и геометрическую картину перехвата. Точки, в которые преследователь
и цель, двигаясь по прямым, приходят одновременно, удовлетворяют условию
$\frac{|AM|}{|BM|} = \frac{u}{v}$. При $u \neq v$ это окружность Аполлония. Перехват
возможен тогда и только тогда, когда прямая, по которой движется цель, пересекает эту
окружность.

%==============================================================================
\subsection{Цель движется по окружности}
%==============================================================================

Если цель сама движется криволинейно, траектория преследователя становится ещё
сложнее. Тем не менее и здесь на многие вопросы можно ответить, если сообразить,
к какому движению погоня выходит со временем.

\begin{example}
Утка плавает по окружности радиуса $R$ с постоянной скоростью $v$. Собака, стартуя
из центра, плывёт со скоростью $u < v$, всё время держа курс на утку. Со временем
движение собаки становится равномерным движением по окружности. Найдите радиус этой
окружности, расстояние между собакой и уткой и ускорение собаки в установившемся
режиме.

\figtikz{images/kin/pur_duck.tikz}{Установившееся движение: собака~$A$ бежит по окружности радиуса $r$, линия визирования $AB$ касается этой окружности.}{fig:pur_duck}

В установившемся режиме взаимное расположение тел не меняется, поэтому собака
движется по концентрической окружности радиуса $r$ с той же угловой скоростью, что и
утка, $\omega = \frac{v}{R}$. Скорость собаки направлена по касательной к её
окружности и одновременно на утку. Значит, линия визирования $AB$ касается
окружности собаки, и угол $OAB$ прямой (рис.~\ref{fig:pur_duck}). Если $\delta$~---
угол между радиусами $OA$ и $OB$, то
\[
r = R\cos\delta,
\qquad
|AB| = R\sin\delta.
\]
С другой стороны, $u = \omega r = \frac{vr}{R}$. Отсюда
\[
r = \frac{u}{v}R,
\qquad
\cos\delta = \frac{u}{v},
\qquad
|AB| = R\sqrt{1 - \frac{u^2}{v^2}}.
\]
Ускорение собаки~--- центростремительное:
\[
a = \frac{u^2}{r} = \frac{uv}{R}.
\]

Проверим ответ формулами линии визирования. Скорость утки перпендикулярна радиусу
$OB$, а линия визирования перпендикулярна радиусу $OA$, поэтому угол между ними равен
углу $\delta$ между радиусами. Тогда
\[
\frac{\Delta r}{\Delta t} = v\cos\delta - u = 0,
\qquad
\Omega = \frac{v\sin\delta}{|AB|} = \frac{v\sin\delta}{R\sin\delta} = \frac{v}{R},
\]
т.\,е. расстояние между телами не меняется, а линия визирования поворачивается
вместе с уткой. Ускорение собаки $u\Omega = \frac{uv}{R}$ совпадает с найденным выше.
\end{example}

Мы не доказывали, что погоня действительно выходит на этот режим; численное
моделирование показывает, что траектория собаки, стартующей из центра, быстро
приближается к найденной окружности. При $u = v$ радиус этой окружности
становится равным $R$, а расстояние~--- нулю: собака неограниченно приближается к
утке, но за конечное время её не догоняет. При $u > v$ собака догоняет утку.

%==============================================================================
\subsection{Как решать задачи преследования}
%==============================================================================

\begin{enumerate}
  \item Нарисуйте линию визирования и отметьте углы между ней и скоростями тел.
  \item Запишите скорость изменения расстояния $\frac{\Delta r}{\Delta t}$ и угловую
  скорость линии визирования $\Omega$. Часто этого уже достаточно: если скорость
  сближения постоянна, время встречи находится делением.
  \item Ищите симметрию. Если конфигурация тел переходит сама в себя при повороте,
  она сохраняется во всё время движения, и задача сводится к одной паре тел.
  \item Если цель движется по прямой, воспользуйтесь инвариантом погони
  $ur + vx = \mathrm{const} - (u^2 - v^2)t$.
  \item Ускорение преследователя с постоянной по модулю скоростью находите как
  $a = u\,|\Omega|$, радиус кривизны~--- как $\rho = \frac{u}{|\Omega|}$.
  \item Для установившихся режимов ищите движение, при котором взаимное расположение
  тел не меняется.
\end{enumerate}

\noindent\textbf{Ключевые понятия:} погоня; линия визирования; скорость сближения;
угловая скорость линии визирования; логарифмическая спираль; инвариант погони;
кривая погони; параллельное сближение; перехват; окружность Аполлония;
установившийся режим погони.

%==============================================================================
\subsection{Задачи}
%==============================================================================

\begin{enumerate}
  \item Шесть жуков сидят в вершинах правильного шестиугольника со стороной $a$ и
  одновременно начинают ползти со скоростью $v$, каждый прямо на соседа по часовой
  стрелке. Через какое время они встретятся? Какой путь проползёт каждый? Каково
  ускорение жука в начальный момент?\\
  \textit{Ответ:} $T = \frac{2a}{v}$, $s = 2a$, $a_0 = \frac{\sqrt{3}\,v^2}{2a}$.

  \item Собака, бегущая со скоростью $2v$, преследует лису, которая бежит по прямой
  со скоростью $v$. В начальный момент расстояние между ними $L$, а скорость лисы
  перпендикулярна линии визирования. Через какое время собака догонит лису? Какие
  пути они пробегут? Каковы ускорение собаки и радиус кривизны её траектории в
  начальный момент?\\
  \textit{Ответ:} $T = \frac{2L}{3v}$; собака пробежит $\frac{4L}{3}$, лиса~---
  $\frac{2L}{3}$; $a_0 = \frac{2v^2}{L}$, $\rho_0 = 2L$.

  \item Три черепахи находятся в вершинах правильного треугольника со стороной $a$ и
  ползут со скоростью $v$, каждая прямо на соседку. Найдите время до встречи и
  радиус кривизны траектории черепахи в момент, когда сторона треугольника
  уменьшилась вдвое.\\
  \textit{Ответ:} $T = \frac{2a}{3v}$, $\rho = \frac{a}{\sqrt{3}}$.

  \item Судно идёт прямым курсом со скоростью $20$~км/ч. Катер находится на
  расстоянии $3$~км от судна, отрезок от судна к катеру составляет с курсом судна угол
  $30^\circ$. С какой наименьшей скоростью должен идти катер, чтобы встретиться с
  судном, и через какое время произойдёт встреча?\\
  \textit{Ответ:} $u_{\min} = 10$~км/ч, $T = \frac{\sqrt{3}}{10}~\text{ч} \approx 10{,}4$~мин.

  \item Утка плавает по окружности радиуса $R$ со скоростью $v$, собака плывёт со
  скоростью $\frac{v}{2}$, всё время держа курс на утку. Найдите радиус окружности, по
  которой в установившемся режиме плывёт собака, расстояние между собакой и уткой и
  ускорение собаки.\\
  \textit{Ответ:} $r = \frac{R}{2}$, $|AB| = \frac{\sqrt{3}}{2}R$, $a = \frac{v^2}{2R}$.
\end{enumerate}

\end{document}
